\begin{ex}
Số liền sau của số 599 là số nào?
\shortans{600}
\loigiai{
Số liền sau 599 là 600.
}
\end{ex}

\begin{ex}
Số liền trước của số 700 là số nào?
\shortans{699}
\loigiai{
Số liền trước 700 là 699.
}
\end{ex}

\begin{ex}
Số gồm 6 trăm, 5 chục và 8 đơn vị là số nào?
\shortans{658}
\loigiai{
Số đó là 658.
}
\end{ex}

\begin{ex}
Số 904 gồm 9 trăm và mấy đơn vị?
\shortans{4}
\loigiai{
Số 904 gồm 9 trăm và 4 đơn vị.
}
\end{ex}

\begin{ex}
Viết số 472 thành tổng: $400 + \dots + 2$. Số cần điền là bao nhiêu?
\shortans{70}
\loigiai{
Số cần điền là 70.
}
\end{ex}

\begin{ex}
Tính: $300 + 500$.
\shortans{800}
\loigiai{
Ta có $300 + 500 = 800$.
}
\end{ex}

\begin{ex}
Tính: $900 - 400$.
\shortans{500}
\loigiai{
Ta có $900 - 400 = 500$.
}
\end{ex}

\begin{ex}
Tính: $250 + 50$.
\shortans{300}
\loigiai{
Ta có $250 + 50 = 300$.
}
\end{ex}

\begin{ex}
Tính: $400 - 80$.
\shortans{320}
\loigiai{
Ta có $400 - 80 = 320$.
}
\end{ex}

\begin{ex}
Trong các số $582, 614, 599, 608$, số nào lớn nhất?
\shortans{614}
\loigiai{
Số lớn nhất là 614.
}
\end{ex}

\begin{ex}
Trong các số $435, 428, 450, 419$, số nào bé nhất?
\shortans{419}
\loigiai{
Số bé nhất là 419.
}
\end{ex}

\begin{ex}
Sắp xếp dãy số $712, 589, 830, 645$ từ bé đến lớn, số đứng ở vị trí thứ ba là số nào?
\shortans{712}
\loigiai{
Thứ tự: 589, 645, 712, 830. Số thứ ba là 712.
}
\end{ex}

\begin{ex}
Điền số thích hợp: $789 < \dots < 791$.
\shortans{790}
\loigiai{
Số ở giữa là 790.
}
\end{ex}

\begin{ex}
Điền số thích hợp: $500 + \dots = 900$.
\shortans{400}
\loigiai{
Số cần điền là: $900 - 500 = 400$.
}
\end{ex}

\begin{ex}
Điền số thích hợp: $1000 - \dots = 300$.
\shortans{700}
\loigiai{
Số cần điền là: $1000 - 300 = 700$.
}
\end{ex}

\begin{ex}
Tính giá trị biểu thức: $400 + 300 - 200$.
\shortans{500}
\loigiai{
Ta có $400 + 300 = 700$, rồi $700 - 200 = 500$.
}
\end{ex}

\begin{ex}
Tính giá trị biểu thức: $800 - 600 + 500$.
\shortans{700}
\loigiai{
Ta có $800 - 600 = 200$, rồi $200 + 500 = 700$.
}
\end{ex}

\begin{ex}
Bác nông dân thu hoạch được 500 bắp ngô nếp và 400 bắp ngô ngọt. Bác thu hoạch được tất cả bao nhiêu bắp ngô?
\shortans{900}
\loigiai{
Bác thu hoạch được: $500 + 400 = 900$ bắp ngô.
}
\end{ex}

\begin{ex}
Một đoàn tàu chở 600 hành khách, đến ga có 200 hành khách xuống tàu. Đoàn tàu còn lại bao nhiêu hành khách?
\shortans{400}
\loigiai{
Đoàn tàu còn lại: $600 - 200 = 400$ hành khách.
}
\end{ex}

\begin{ex}
Kho gạo thứ nhất có 350 bao gạo, kho thứ hai có nhiều hơn kho thứ nhất 50 bao gạo. Kho thứ hai có bao nhiêu bao gạo?
\shortans{400}
\loigiai{
Kho thứ hai có: $350 + 50 = 400$ bao gạo.
}
\end{ex}

\begin{ex}
Tìm số lớn nhất có ba chữ số mà chữ số hàng trăm là 7.
\shortans{799}
\loigiai{
Số lớn nhất là 799.
}
\end{ex}

\begin{ex}
Tìm số bé nhất có ba chữ số mà chữ số hàng trăm là 4.
\shortans{400}
\loigiai{
Số bé nhất là 400.
}
\end{ex}

\begin{ex}
Số 1000 trừ đi số lớn nhất có ba chữ số thì được kết quả bằng bao nhiêu?
\shortans{1}
\loigiai{
Số lớn nhất có 3 chữ số là 999. Hiệu: $1000 - 999 = 1$.
}
\end{ex}

\begin{ex}
Có bao nhiêu số tròn trăm từ 200 đến 800?
\shortans{7}
\loigiai{
Các số: 200, 300, 400, 500, 600, 700, 800. Có 7 số.
}
\end{ex}

\begin{ex}
Tính tổng các chữ số của số bé nhất có ba chữ số khác nhau (102).
\shortans{3}
\loigiai{
Số 102 có tổng các chữ số: $1 + 0 + 2 = 3$.
}
\end{ex}
